% !TeX root = ../main.tex

\section{Introduction}

% Modern software systems are increasingly susceptible to exploitation and compromise through internal defects. 


In recent years, the open-source software (OSS) ecosystem has witnessed substantial growth \cite{octoverse,sonatype}, providing developers with a vast array of options for using free software. It has become a common practice for developers to adopt third-party OSS components to accelerate development and reduce costs \cite{koschke2016software}. Consequently, the vulnerability landscape has expanded beyond traditional proprietary software to encompass open-source software (OSS). The increasing vulnerabilities in open-source software have raised significant concerns, particularly regarding supply chain attacks \cite{crowdstrike}. Attackers can exploit these vulnerabilities, affecting not only the open-source software itself but also its downstream users, thus amplifying potential financial gains and cyber chaos \cite{arcticwolf}.


Notably, downstream software shares or reuse logic from upstream OSS components. As these OSS components evolve and downstream software undergoes customization, the code initially derived from OSS components diverges from the original one at both ends. Consequently, it leads to challenges in mitigating \textit{Vulnerable Code Clones} (VCCs) \cite{pham2010detection}. When facing incompatible changes or customized OSS components, mitigating VCCs by replacing newer OSS components becomes infeasible or cumbersome, especially in languages like C/C++ that lacks a dominant package manager. 


In this context, developers often resort to manually finding and backporting security patches into downstream software. However, it is time-consuming and resource-intensive. Notably, 80\% of attacks leveraged the vulnerabilities reported more than three years ago \cite{businessinsights}, indicating that the issue of VCCs has not been~fully~resolved.


Typically, developers employ vulnerable code clone (VCC) discovery techniques \cite{kim2017vuddy, xiao2020mvp, woo2022movery,woo2023v1scan} to identify reused components (such as files, functions, or fragments) in a target project program. They can also use learning-based methods for classifying vulnerable functions \cite{russell2018automated, church2017word2vec, wang2023graphspd, zhou2019devign, cui2020vuldetector}. However, the effectiveness of vulnerability detection is often hindered in cases of \textit{Vulnerabilities with Multiple Fixing Functions} (VM), where fixing patches span across multiple functions or files. Unfortunately, VM are inevitably prevalent in real world due to complex inter-procedural program logic, posing challenges for accurate VM detection.

\textbf{Limitations of Existing Approaches.} Existing research often addresses the detection of \textit{Vulnerabilities with Multiple Fixing Functions} (VM) by dividing the task into matching vulnerability characteristics within individual functions (\textit{i.e.,} matching-one-in-all approach). Clone-based techniques \cite{kim2017vuddy, xiao2020mvp, woo2022movery, woo2023v1scan} identify vulnerable code clones (VCCs) within functions, while learning-based approaches represent functions in abstract spaces and classify them as vulnerable or not \cite{russell2018automated, church2017word2vec, wang2023graphspd, zhou2019devign, cui2020vuldetector}. However, these approaches assume equal importance for all fixing functions in a vulnerability, potentially leading to over-representation resulting in false alarms. Additionally, while software composition analysis (SCA) techniques \cite{woo2021centris, yang2022modx, jiang2023third, wu2023ossfp, woo2023v1scan} can detect VM, they are primarily designed for identifying reused OSS components, which may overlook specific vulnerability or fixing characteristics, resulting in false alarms.


\textbf{Challenge.} Naturally thinking, assigning all fixing equal importance is straightforward and may not accurately represent the actual characteristics of vulnerabilities. Some fixing functions may carry adequate information to represent a vulnerability, while others may not. However, it is challenging to measure and weigh the unique fingerprints of each fixing function to a vulnerability. The primary challenge lies in how to capture and distinguish the uniqueness of fixing functions and identify the critical ones for a vulnerability.


\textbf{Our Approach.} To tackle this challenge, we introduce \tool, a novel approach for \underline{V}ulnerabilities-with-\underline{Mu}ltiple-Fixing-Function-\underline{D}etection. Unlike traditional methods that treat all fixing functions equally, \tool prioritizes the fixing functions and selects the most critical functions for vulnerability representation (see Section \ref{sec:cfs}). It involves matching the signatures \cite{xiao2020mvp, kim2017vuddy} of these critical functions. To cope with the potential decrease in recall due to excluding the remaining fixing functions, we employ \todo{semantic equivalent statement matching}, which consists of program rephrasing \cite{visser2001survey} (see Section \ref{sec:rephrasing}) and \todo{contextual semantic equivalent statement mapping} (see Section \ref{sec:stage2:matching}). The program rephrasing aims to uncover more VM by creating another signature for the rephrased form of the critical functions, while \todo{contextual semantic equivalent statement mapping} aims to match the composing statements precisely based on their contextual semantic~equivalence.

\minor{The insight of the contextual semantic equivalent statement mapping stems from our observation that some syntactically identical statements serve different contextual purposes. Existing approaches add new context similarity constraints using syntactically equivalent statements, which does not correct the error from the essential. We claim that ``statement equivalence'' should be redefined not only from syntactic similarity, but also from their contextual similarity. To our best knowledge, \tool is the first to discriminate critical functions from fixing functions and consider both program rephrasing and contextual semantic equivalent statements to identify vulnerabilities with high quantity and quality.}

\textbf{Evaluation.} We implemented \tool and generated vulnerability signatures from \todo{810} CVEs. We conducted an evaluation using a dataset comprised of the top \todo{1,000} C/C++ real-world projects selected from GitHub. Comparing \tool with state-of-the-art techniques, % including \todo{V1Scan}, \todo{Vuddy}, \todo{MVP}, and \todo{Movery}.
\todo{\tool achieves a f1-score of \todo{0.84}, outperforming the best state-of-the-art at \todo{17.6\%}}.
\tool discovered \todo{275} VM across \todo{104} projects, with \todo{42} of them confirmed by developers and \todo{5} of them assigned with corresponding new CVE identifiers. We observed the robustness of \tool regarding the ixing function number in the VM. We conducted an ablation study to observe the contribution of key components in \tool. We also measured the threshold sensitivity and performance of \tool. The results demonstrate~\tool's effectiveness and efficiency in identifying VM within real-world~projects. 

\textbf{Contribution.} We summarize our contributions as follows:
\begin{itemize}[leftmargin=*]
 \item We propose \tool, a novel approach to detect Vulnerabilities with Multiple Fixing Functions (VM) in real-world projects.
 \item We evaluate \tool on our ground truth collected from real-world C/C++ projects and demonstrate its effectiveness compared with the state-of-the-art approaches.
 \item \tool has discovered \todo{275} vulnerabilities in \todo{84} C/C++ projects. \todo{42} are confirmed by developers, \todo{5} has been assigned with CVE identifiers, demonstrating its effectiveness in real-world projects.
\end{itemize}






